\section{Observable upper bound}\label{sec:direct-upper}
We construct an estimator at the quantile split from
Section~\ref{sec:model}. The classification and allocation lemmas
used below are proved in Appendix~\ref{sec:uf-foundations}.

\begin{theorem}[Upper bound at the fixed quantile split]\label{thm:direct-upper}
For every admissible $(\mu,\rho,B)$, there is an observable estimator
with uniform expected work at most $B$ and worst RMS error at most
$30001\Q(B)$.
\end{theorem}

If $\Q(B)=\rho$, use the constant estimator $\rho/2$.
Otherwise $h<t$, because $h\ge t$ would imply $T_h=\rho$.
Fix
\[
 \sigma=\kappa,\qquad \tau=\nu|_{[h,1/2]},\qquad
 \zeta=\nu|_{(0,h)},\qquad z=\zeta((0,1/2]),\quad
 H=\int d^{-1}\,\tau(dd).
\]
Then $\mu=\sigma+\tau+\zeta$ and $T_h=z+hH$.
These are analytical pieces only. The estimator never observes which
piece produced its genuine coin.
All likelihood evaluations and design parameters use the known law.
The complexity measure counts physical coin observations, including those
given to either side of the race; it does not assert a bound on arithmetic
time, memory, or random-number generation.

\subsection{An uncapped reference bank with a capped observable race}

Use the fixed decomposition above, put
$z=\zeta((0,1/2])$, and assume $\Gamma_\sigma,H_\tau<\infty$.
It suffices to treat $z\le\rho/128$: otherwise the constant $\rho/2$
has risk at most $64z$. For a finite submeasure $\xi$, define explicitly
\[
 a_j(\xi)=\int_{I_j}d^{-2}\,\xi(dd),\qquad
 A_\xi=\sum_j a_j(\xi),\qquad H_\xi=\int d^{-1}\,\xi(dd),
\]
\[
 \Gamma_\xi=A_\xi+\sum_j a_j(\xi)\log(A_\xi/a_j(\xi)).
\]
Zero coordinates are omitted, and all three quantities are zero for the
zero measure. These are exact inverse-square costs, not the rounded
canonical costs used in the lower bound. Write $A=A_\sigma$, $H=H_\tau$.
If $\sigma\ne0$, use all its occupied bands, with weights
$\pi_j=a_j(\sigma)/A$ and references $b_j/2$ and $b_j/4$.
For an internal allowance $u<\rho$, set $\epsilon=u/\rho$.
Stop negative if a negative reference likelihood against fair reaches
$1/(\epsilon\pi_j)$; stop positive if
\begin{equation}\label{eq:split-bank-ratio}
 \frac{\sum_j\pi_jP_{1/2+b_j/4,n}}
      {\int P_{1/2-d,n}\,\mu(dd)}\ge1/u.
\end{equation}
Give negative crossings priority. A zero denominator triggers a positive
stop. The known all-negative mixture in the denominator is the full law
$\mu$, regardless of the decomposition.

By Lemma~\ref{lem:uf-mixture-bank}, a wrong negative crossing has
conditional positive-sign probability at most $\epsilon$, and a positive
crossing has all-negative-mixture probability at most $u$. Thus the total
probability of a wrong bank stop is at most $2u$ over $\mathcal Q_\rho$.
Termination outside $\sigma$ is not assumed.

Race this bank against the ordinary classifier of Lemma~\ref{lem:uf-ordinary}, using
alternating observations of the same physical coin. At an ordinary cap
$M$, its conditional work certificate is
\[
 o_M(d)=\min\{M,17d^{-2}[1+\log_+(Md^2)]\}.
\]
For any finite measure $\xi$, put
\[
 \eta_\xi(M)=\int\min\{1,(Md^2)^{-2}\}\,\xi(dd),\qquad
 O_\xi(M)=\int o_M(d)\,\xi(dd).
\]
The ordinary error under $\mu$ is at most $\eta_\mu(M)$.
Stop as soon as either component stops and return that output.
On completed counterfactual streams, the physical work is at most
\[
 2\min\{T_{\rm bank},T_{\rm ordinary}\}\le2M.
\]
A wrong output implies either a wrong bank stop or an ordinary error.
When $\sigma=0$, omit the bank and use the ordinary classifier alone.

On a $\sigma$-almost-every gap $d$ in band $j$, the matching negative
reference has drift at least $4d^2/3$ and log overshoot below one.
Its expected crossing time is at most
\begin{equation}\label{eq:split-negative-work}
 \frac3{4d^2}[1+\log(1/\epsilon)+\log(1/\pi_j)].
\end{equation}
On a positive gap, the event $\widehat p_n-1/2\ge d/2$ gives
\begin{equation}\label{eq:split-null-laplace}
 \frac{P_{-,\mu,n}}{P_{{\rm f},n}}
 \le\int e^{-2nx^2}\,\mu(dx)
 \le z+\frac A{2en}+\frac H{2\sqrt{en}}.
\end{equation}
The last two bounds use, respectively,
$\sup_xx^2e^{-2nx^2}=1/(2en)$ and
$\sup_xxe^{-2nx^2}=1/(2\sqrt{en})$.
For $n\ge d^{-2}$ the denominator is at most $z+Ad^2+Hd$,
while the matching positive reference contributes at least
$\pi_je^{nd^2/3}$ relative to fair. Equation~\eqref{eq:uf-bank-positive} therefore bounds positive expected bank work by
\begin{equation}\label{eq:split-positive-work}
 \frac6{d^2}\left[
 1+\log(1/\epsilon)+\log(1/\pi_j)
 +\log_+\frac{z+Ad^2+Hd}{\rho}\right].
\end{equation}
These arguments prove finite conditional bank work on $\sigma$.
On $\tau+\zeta$, the ordinary cap guarantees termination of the race.

Since $z/\rho\le1$,
\[
 \log_+\frac{z+Ad^2+Hd}{\rho}
 \le\log3+\log_+(Ad^2/\rho)+\log_+(Hd/\rho).
\]
Using $\sup_{x\ge1}\log x/x=1/e$ and
$\sup_{x\ge1}\log x/x^2=1/(2e)$ gives
\begin{equation}\label{eq:split-surcharge}
 \int q(d)d^{-2}\log_+\frac{z+Ad^2+Hd}{\rho}\,\sigma(dd)
 \le(\log3+1/e)A+\frac{H^2}{2e\rho}.
\end{equation}
Here $\int q\,d\sigma\le\rho$ and $q\le1$ follow from the original
global rarity class; no separate conditional rarity assumption is used.
Integrating \eqref{eq:split-negative-work} and
\eqref{eq:split-positive-work} therefore bounds the bank work on $\sigma$ by
\begin{equation}\label{eq:split-bank-work}
 6\Gamma_\sigma+6A\log(1/\epsilon)
       +6(\log3+1/e)A+\frac{3H^2}{e\rho}.
\end{equation}

\subsection{A resolved-law cap schedule}

Let $\chi=\sigma+\tau$ and set
\[
 \delta_\ell=\rho16^{-(\ell+1)},\qquad
 u_\ell=\delta_\ell/4,\qquad \epsilon_\ell=u_\ell/\rho.
\]
Choose $M_\ell$ as the least power of four, at least four, such that
$\eta_\chi(M_\ell)\le\delta_\ell/8$.
Every chosen cap has a predecessor because $\eta_\chi(4)=1-z$.
The inequality $\eta_\chi(4M)\ge\eta_\chi(M)/16$ implies
\begin{equation}\label{eq:split-cap-selection}
 \delta_\ell/128<\eta_\chi(M_\ell)\le\delta_\ell/8.
\end{equation}
Thus the selected powers of four are distinct and increasing.

If $z>0$, let $J$ be the largest integer with $\delta_J\ge8z$.
Then
\begin{equation}\label{eq:split-terminal-level}
 8z\le\delta_J<128z.
\end{equation}
For $\ell\le J$,
$\eta_\mu(M_\ell)\le\eta_\chi(M_\ell)+z\le\delta_\ell/4=u_\ell$.
The raced classifier has total error at most $3u_\ell\le\delta_\ell$
and false-positive probability at most $2u_\ell\le\delta_\ell$.
When $z=0$, these conclusions hold at every level.

Since $\log(1/\epsilon_\ell)=\log64+\ell\log16$ and
$A\le\Gamma_\sigma$, \eqref{eq:split-bank-work} is at most
$40(\ell+1)\Gamma_\sigma+2H^2/\rho$.
A uniform raced-work certificate is consequently
\begin{equation}\label{eq:split-level-work}
 C_\ell=80(\ell+1)\Gamma_\sigma+4H^2/\rho
                  +2O_\tau(M_\ell)+2zM_\ell.
\end{equation}
It remains a valid, possibly loose, bound when $\sigma=0$ and the bank is
omitted.

The binary level-zero output has second moment at most
$\rho+\delta_0=17\delta_0$. An adjacent difference, run on successive
unused blocks of the same coin, has second moment at most
$\delta_\ell+\delta_{\ell-1}=17\delta_\ell$.
Set $v_\ell=17\delta_\ell$, $k_0=C_0$,
$k_\ell=C_\ell+C_{\ell-1}$ for $\ell\ge1$, and
$S_J=\sum_{\ell\le J}\sqrt{v_\ell k_\ell}$.
Independent-Poisson allocation across these corrections
(Lemma~\ref{lem:uf-poisson})
gives RMS at most
\begin{equation}\label{eq:split-poisson}
 \sqrt{\delta_J^2+S_J^2/B}.
\end{equation}
Every such finite-level procedure has finite deterministic local caps.
In the finite-law theorem the reference bank is finite; its stopping
events are measurable and its positive reference is a finite probability
mixture. No countable-support extension is needed here.

\subsection{Controlling mixed work without an extra logarithm}

For every finite $\tau$, the following bounds hold:
\begin{equation}\label{eq:split-elementary-kernels}
 O_\tau(M)\le21H\sqrt M,\qquad
 \eta_\sigma(M)\le A/M,\qquad
 \eta_\tau(M)\le H/\sqrt M.
\end{equation}
The latter two are pointwise inequalities before integration.
For the first, put $x=\sqrt M\,d$. The ratio of $o_M(d)$ to
$\sqrt M/d$ is at most one for $x\le1$, and at most
$17(1+2\log x)/x\le34/\sqrt e<21$ for $x\ge1$.
Lemma~\ref{lem:uf-kernel} supplies the homogeneous ordinary kernel bound:
\begin{equation}\label{eq:split-ordinary-kernel}
 \sum_{k\ge1}\sqrt{\eta_\tau(4^k)O_\tau(4^k)}\le 208H.
\end{equation}
Its proof applies without normalization of $\tau$: dyadic binning yields
the summable kernels $2^{k-j}$ for $j\ge k$, and
$2^{3(j-k)}$, $2^{j-k}(1+k-j)$ for $j<k$.
Apply $\sqrt{xy}\le(x+y)/2$ and sum these kernels.

Split levels according to which term of
$\eta_\chi=\eta_\sigma+\eta_\tau$ is at least half the total.
On levels where $\eta_\tau$ dominates,
$\delta_\ell\le256\eta_\tau(M_\ell)$, so distinctness of the caps and
\eqref{eq:split-ordinary-kernel} give a total
$\sum\sqrt{\delta_\ell O_\tau(M_\ell)}\le3328H$.
On the other levels,
\[
 \delta_\ell\le256\eta_\sigma(M_\ell)\le256A/M_\ell,
 \qquad M_\ell\le256A/\delta_\ell.
\]
Therefore
\[
 \sqrt{\delta_\ell O_\tau(M_\ell)}
 \le4\sqrt{21}\,H^{1/2}A^{1/4}\delta_\ell^{1/4}.
\]
The sum of $\delta_\ell^{1/4}$ over all levels is $\rho^{1/4}$.
The arithmetic--geometric mean inequality proves
\begin{equation}\label{eq:split-mixed-sum}
 \sum_{\ell\le J}\sqrt{\delta_\ell O_\tau(M_\ell)}
 \le C_{\rm mix}[H+\sqrt{\rho A}].
\end{equation}

The work on the unresolved mass is controlled separately.
On levels where $\eta_\sigma$ dominates,
\begin{equation}\label{eq:split-unresolved-sigma}
 \sum\sqrt{z\delta_\ell M_\ell}
 \le16(J+1)\sqrt{zA}\le\sqrt{2\rho A}.
\end{equation}
The last inequality uses $z\le\delta_J/8$ and
$(J+1)4^{-(J+1)}\le1/4$.
On levels where $\eta_\tau$ dominates,
$\sqrt{M_\ell}\le256H/\delta_\ell$, so
\begin{equation}\label{eq:split-unresolved-tau}
 \sum\sqrt{z\delta_\ell M_\ell}
 \le256H\sqrt z\sum_{\ell=0}^J\delta_\ell^{-1/2}
 \le\frac{1024}{3}H\sqrt{z/\delta_J}
 \le\frac{1024}{3\sqrt8}H.
\end{equation}
When $z=0$ these contributions vanish.

\begin{proof}[Completion of Theorem~\ref{thm:direct-upper}]
The shifted adjacent-level sum has a factor $1/4$, so
\[
 S_J\le\frac{5\sqrt{17}}4
                \sum_{\ell=0}^J\sqrt{\delta_\ell C_\ell}.
\]
Use \eqref{eq:split-level-work}, \eqref{eq:split-mixed-sum},
\eqref{eq:split-unresolved-sigma}, \eqref{eq:split-unresolved-tau},
and $A\le\Gamma_\sigma$. The elementary sums
$\sum_{\ell\ge0}4^{-\ell}\sqrt{\ell+1}\le16/9$ and
$\sum_{\ell\ge0}4^{-\ell}=4/3$ then give
\[
 S_J\le C_{\rm sum}[\sqrt{\rho\Gamma_\sigma}+H].
\]
One valid explicit choice is $C_{\rm sum}=30000$. Indeed, writing
$R=\sqrt{\rho\Gamma_\sigma}$, the component bounds above give
\[
 \sum_{\ell\le J}\sqrt{\delta_\ell C_\ell}
 \le21R+\frac{15535}{3}H.
\]
Here the core term is at most $4R$, the fixed tail surcharge is $2H/3$,
the ordinary mixed term is at most $15R+5007H$, and the unresolved
term is at most $2R+512H/3$. Use $\sqrt{21}<5$, $\sqrt2<3/2$ and
$5\sqrt{17}/4<165/32$ to obtain
\[
 S_J\le\frac{3465}{32}R+\frac{854425}{32}H
       \le30000(R+H).
\]
For $z>0$, \eqref{eq:split-terminal-level} and
\eqref{eq:split-poisson} imply
$R(B)\le128z+C_{\rm sum}[\sqrt{\rho\Gamma_\sigma}+H]/\sqrt B$.
For $z=0$, put $D=\sqrt{\rho\Gamma_\sigma}+H>0$ and choose the least
finite $J\ge0$ such that $\delta_J\le D/\sqrt B$. The bound on $S_J$
holds at this finite terminal level. Equation~\eqref{eq:split-poisson}
therefore gives RMS at most
\[
 \sqrt{1+30000^2}\,D/\sqrt B<30001D/\sqrt B.
\]
Thus this branch has an actual finite-level estimator; no limiting
or infinite-level construction is needed.
The constant estimator covers $z>\rho/128$.
For the fixed split, $z+hH=T_h$ and $\sigma=\kappa$.
The preceding bound and the constant estimator therefore give
$R(B)\le C\min\{\rho,T_h+h\sqrt{\rho\Gamma_\kappa}\}=C\Q(B)$.
The preceding constants permit $C=30001$ for this upper bound; no
optimality of that constant is claimed. The number of levels and every
local cap are finite, while the total Poisson replication count is
unbounded but almost surely finite and has the stated expected budget.
Since the law is finite, the reference bank is finite throughout.
\end{proof}
